% META: {"id": "TSPE-DERCONV-CO-002", "chapitre": "TSPE-DERIVATION-CONVEXITE", "type_objet": "corrige", "exercice_ref": "TSPE-DERCONV-EX-002", "capacites_codes": ["C1"], "sources_inspiration": [], "status": "approved"}
% BEGIN-VERIFY
% from sympy import *
% x = symbols('x')
% u = x**2 + 3*x + 2
% du = diff(u, x)
% f = ln(u)
% fp = diff(f, x)
% assert simplify(fp) == (2*x + 3)/(x**2 + 3*x + 2)
% # domaine: u > 0, i.e. x < -2 ou x > -1
% assert u.subs(x, 0) == 2
% assert u.subs(x, -3) == 2
% END-VERIFY
\begin{corrige}{TSPE-DERCONV-EX-002}

\textbf{Question 1.} $x^2+3x+2 = (x+1)(x+2) > 0$ si et seulement si $x \in \left]-\infty, -2\right[ \cup \left]-1, +\infty\right[$.

Donc $D_f = \left]-\infty, -2\right[ \cup \left]-1, +\infty\right[$.

\commentaireMarge{Factoriser le polynome sous le logarithme.}

\textbf{Question 2.} On pose $u(x) = x^2+3x+2$. Alors $u'(x) = 2x+3$.

\[ f'(x) = \frac{u'(x)}{u(x)} = \frac{2x+3}{x^2+3x+2}. \]

\end{corrige}
